\chapter{The flag algebra $\mathcal{A}^\sigma$}
\label{chap:algebra}

Formal real combinations of flags form a free module; a product defined through
pair densities and a quotient by Razborov's averaging relations turn it into a
commutative $\mathbb{R}$-algebra, the flag algebra $\mathcal{A}^\sigma$.
(Source: \texttt{FlagAlgebra/FlagAlgebra.lean}.)

\begin{definition}[Size-tagged flag]
  \label{def:fin-flag}
  \lean{FlagAlgebras.FinFlag}
  \uses{def:flag-with-size}
  \leanok
  A \emph{size-tagged flag} is a dependent pair $\langle n, F\rangle$ of a vertex
  count $n$ and a flag $F$ on $\mathrm{Fin}\,n$. These are the basis indices of
  the algebra.
\end{definition}

\begin{definition}[Flag vector]
  \label{def:flag-vector}
  \lean{FlagAlgebras.FlagVector}
  \uses{def:fin-flag}
  \leanok
  A \emph{flag vector} is a finitely supported real combination of size-tagged
  flags, i.e.\ an element of $\mathrm{FinFlag}\,\sigma \to_{0} \mathbb{R}$: the
  underlying module before quotienting.
\end{definition}

\begin{definition}[Basis vector]
  \label{def:basis-vector}
  \lean{FlagAlgebras.basisVector}
  \uses{def:flag-vector}
  \leanok
  $\mathrm{basisVector}\,F$ is the unit vector supported on the single flag $F$;
  every construction reduces to these.
\end{definition}

\begin{definition}[Flag multiplication]
  \label{def:flag-mul}
  \lean{FlagAlgebras.flagMul, FlagAlgebras.flagMulWithSize}
  \uses{def:basis-vector, def:flag-density-n, def:flag-with-size}
  \leanok
  The product of two flags is expanded onto flags $G$ of a target size $\ell$,
  weighted by the pair density:
  \[ F \cdot F' \;=\; \sum_{G} d_2(F,F';G)\,\cdot G, \]
  taken at the minimal size $\ell = |F| + |F'| - |\sigma|$
  ($\mathrm{flagMul}$); the choice of $\ell$ is irrelevant modulo the averaging
  relations.
\end{definition}

\begin{definition}[Zero space]
  \label{def:zero-space}
  \lean{FlagAlgebras.ZeroSpace}
  \uses{def:zero-set, def:flag-vector}
  \leanok
  The \emph{zero space} $\mathrm{ZeroSpace}\,\sigma$ is the $\mathbb{R}$-submodule
  spanned by all of Razborov's averaging relations (the set
  $\mathrm{zeroSet}\,\sigma$). Quotienting by it identifies each flag with its
  density expansion.
\end{definition}

\begin{definition}[The flag algebra]
  \label{def:flag-algebra}
  \lean{FlagAlgebras.FlagAlgebra}
  \uses{def:flag-vector, def:zero-space, def:flag-vector-eqv}
  \leanok
  $\mathcal{A}^\sigma = \mathrm{FlagVector}\,\sigma / {\sim_v}$, flag vectors
  modulo the zero space. It carries a commutative $\mathbb{R}$-algebra structure
  (built below), with quotient map written $[\![\,\cdot\,]\!]$.
\end{definition}

\begin{theorem}[The zero space is an ideal]
  \label{thm:zero-space-ideal}
  \lean{FlagAlgebras.flagVector_mul_zeroSpace}
  \uses{def:zero-space, def:flag-mul}
  \leanok
  Multiplying any flag vector by an averaging relation gives another element of
  the zero space. Hence multiplication descends to a well-defined operation on
  the quotient.
\end{theorem}
\begin{proof}\leanok
  Reduce to a basis vector times a generator $\mathrm{zeroElement}\,F\,\ell$ and
  apply the chain rule (Theorem~\ref{thm:density-chain-rules}), which is exactly
  the statement that the product of the relation is again a telescoping sum of
  relations.
\end{proof}

\begin{theorem}[Associativity modulo $\sim_v$]
  \label{thm:mul-assoc}
  \lean{FlagAlgebras.flagVector_mul_assoc}
  \uses{def:flag-mul, def:flag-density-n}
  \leanok
  $(f\cdot g)\cdot h \sim_v f\cdot(g\cdot h)$. Associativity fails on the nose on
  flag vectors but holds modulo the zero space, so it lifts to genuine
  associativity on $\mathcal{A}^\sigma$.
\end{theorem}
\begin{proof}\leanok
  \uses{thm:triple-mul-density}
  Both sides expand, via the triple-density identity
  (Theorem~\ref{thm:triple-mul-density}), to the same sum
  $\sum_{G} d_3(F_1,F_2,F_3;G)\cdot G$.
\end{proof}

\begin{theorem}[Product equals density expansion]
  \label{thm:quot-mul-density}
  \lean{FlagAlgebras.basisVector_quot_mul_eq_flagMul_quot}
  \uses{def:flag-algebra, def:flag-mul}
  \leanok
  In $\mathcal{A}^\sigma$, $[\![F]\!]\cdot[\![G]\!] = [\![\,\mathrm{flagMul}\,F\,G\,]\!]$:
  the ring product of two flags is the class of their density expansion.
\end{theorem}
\begin{proof}\leanok
  Unfold the multiplication on the quotient, which is defined by
  $\mathrm{flagMul}$ and well-defined by Theorem~\ref{thm:zero-space-ideal}.
\end{proof}

\begin{theorem}[Same-size flags sum to one]
  \label{thm:sum-flags-eq-one}
  \lean{FlagAlgebras.sum_flagWithSize_eq_one}
  \uses{def:flag-algebra, def:basis-vector}
  \leanok
  For every $\ell \ge |\sigma|$, $\sum_{F : |F|=\ell} [\![F]\!] = 1$ in
  $\mathcal{A}^\sigma$: the sum over all flags of a fixed size is the unit. This
  normalization is used pervasively in density proofs.
\end{theorem}
\begin{proof}\leanok
  \uses{thm:basis-quot-eq-sum}
  Specialize the density-expansion identity
  (Theorem~\ref{thm:basis-quot-eq-sum}) to the empty flag, whose density in every
  size-$\ell$ flag is $1$.
\end{proof}

\section{Supporting constructions}

\begin{definition}[Flag of fixed size]
  \label{def:flag-with-size}
  \lean{FlagAlgebras.FlagWithSize}
  \uses{def:flag}
  \leanok
  $\mathrm{FlagWithSize}\,\sigma\,n = \mathrm{Flag}\,\sigma\,(\mathrm{Fin}\,n)$:
  a flag on the fixed $n$-element carrier.
\end{definition}

\begin{definition}[Density expansion]
  \label{def:flag-expansion}
  \lean{FlagAlgebras.flagExpansion}
  \uses{def:basis-vector, def:flag-density-n, def:flag-with-size}
  \leanok
  $\mathrm{flagExpansion}\,F\,\ell = \sum_{F'} d_1(F;F')\,\cdot F'$, the averaged
  expansion of $F$ onto size-$\ell$ flags.
\end{definition}

\begin{definition}[Averaging relation]
  \label{def:zero-element}
  \lean{FlagAlgebras.zeroElement}
  \uses{def:basis-vector, def:flag-expansion}
  \leanok
  Razborov's generating relation
  $\mathrm{zeroElement}\,F\,\ell = F - \mathrm{flagExpansion}\,F\,\ell$,
  identified with $0$ in the algebra.
\end{definition}

\begin{definition}[Relation set]
  \label{def:zero-set}
  \lean{FlagAlgebras.zeroSet}
  \uses{def:zero-element}
  \leanok
  The set of all averaging relations $\mathrm{zeroElement}\,F\,\ell$ (with
  $|F| \le \ell$); its span is the zero space.
\end{definition}

\begin{definition}[Flag equivalence of vectors]
  \label{def:flag-vector-eqv}
  \lean{FlagAlgebras.flagVectorEqv}
  \uses{def:zero-space}
  \leanok
  $f \sim_v g$ iff $f - g \in \mathrm{ZeroSpace}\,\sigma$: the two vectors
  represent the same flag-algebra element.
\end{definition}

\begin{theorem}[Quotient map is a homomorphism]
  \label{thm:quot-hom}
  \lean{FlagAlgebras.add_quot, FlagAlgebras.mul_quot, FlagAlgebras.smul_quot, FlagAlgebras.sum_quot}
  \uses{def:flag-vector, def:flag-algebra}
  \leanok
  The quotient map $[\![\,\cdot\,]\!]$ commutes with addition, multiplication,
  scalar multiplication, and finite sums.
\end{theorem}

\begin{theorem}[Triple product expansion]
  \label{thm:triple-mul-density}
  \lean{FlagAlgebras.three_flag_mul_eqv_sum_tripleDensity}
  \uses{def:basis-vector, def:flag-vector-eqv, def:flag-with-size, def:flag-density-n, def:flag-mul, thm:density-chain-rules}
  \leanok
  The product of three flags is, modulo $\sim_v$, the size-$\ell$ sum weighted by
  the triple density $d_3$. This is the engine behind associativity
  (Theorem~\ref{thm:mul-assoc}).
\end{theorem}

\begin{theorem}[Flag class as density expansion]
  \label{thm:basis-quot-eq-sum}
  \lean{FlagAlgebras.basisVector_quot_eq_sum}
  \uses{def:basis-vector, def:flag-with-size, def:flag-density-n, def:flag-algebra, thm:quot-hom, def:flag-expansion}
  \leanok
  In the algebra, $[\![F]\!] = \sum_{F'} d_1(F;F')\,[\![F']\!]$ over size-$\ell$
  flags: the averaging relation written as a $[\![\,\cdot\,]\!]$ identity.
\end{theorem}

\begin{theorem}[Nontriviality]
  \label{thm:flag-algebra-nontrivial}
  \lean{FlagAlgebras.flagAlgebra_zero_ne_one}
  \uses{def:flag-algebra, def:zero-space, def:basis-vector, def:flag-density-n, thm:density-chain-rules}
  \leanok
  $0 \ne 1$ in $\mathcal{A}^\sigma$: the flag algebra is nontrivial (so its
  $\mathrm{CommRing}$ structure is not the zero ring).
\end{theorem}
\begin{proof}\leanok
  If $1$ lay in $\mathrm{ZeroSpace}\,\sigma$ it would be a finite combination of
  averaging relations. But the linear functional $\varphi(G') = d_1(G';F)$
  (density into a fixed large flag $F$) sends every generator
  $\mathrm{zeroElement}\,G\,\ell$ to $0$ by the chain rule
  (Theorem~\ref{thm:density-chain-rules}), while $\varphi(1) = 1$ --- a
  contradiction.
\end{proof}
